% !TeX root = ../../Book1.tex
\section{淡收敛与弱收敛}
\label{b1:sec:1201}

把一个测度看成积分泛函以后，测度列的收敛便可以通过积分值来定义。
不过，测试函数的范围不能省略。同一列概率测度，用紧支集函数测试可能趋于零，
用常函数测试却始终得到一。这里的差别不是记号上的习惯，而是是否保留整体质量。

\ProseParagraphBreak

本节先在局部紧 Hausdorff 空间上接续前一章的表示定理，再转到度量空间上的概率测度。
本章未特别说明时，测度均为正 Borel 测度；需要符号或复测度时会另行指出。

\subsection{测试函数}
\label{b1:subsec:120101}

对拓扑空间 \(X\)，记
\[
 C_b(X)\coloneqq\{\,f:X\to\mathbb R:f\text{ 连续且有界}\,\},
 \quad \|f\|_\infty\coloneqq\sup_{x\in X}|f(x)|.
\]
\symidx{\(C_b(X)\)：有界连续函数空间}
它在一致范数下是 Banach 空间。一致 Cauchy 列的逐点极限仍有界，且收敛一致，
因而极限连续。以下用实值函数测试已经足够；若要使用复值函数，分别测试实部和虚部即可。

\ProseParagraphBreak

当 \(X\) 局部紧 Hausdorff 时，有
\[
 C_c(X)\subset C_0(X)\subset C_b(X).
\]
三个空间分别允许我们观察紧区域、无穷远处逐渐减弱的区域，以及整个空间。
若 \(X\) 紧，三者相同；若 \(X\) 非紧，常函数 \(1\) 属于 \(C_b(X)\)，
却不属于 \(C_0(X)\)。对局部有限的测度，\(C_c(X)\) 上的积分总是有限的；
要让全部有界测试函数都可积，则必须要求测度有限。

\ProseParagraphBreak

有限符号或复 Radon 测度与 \(C_0(X)^*\) 的等距对应已在%
\cref{b1:thm:riesz-markov-c0}中建立。它给出了一种现成的弱*\WeakStarJoin{}收敛：
对每个 \(f\in C_0(X)\)，积分 \(\int f\,\mathrm{d}\mu_n\) 收敛。
但本章还要使用 \(C_c(X)\) 和 \(C_b(X)\)。每当比较两种收敛时，
应先写出测试空间，再判断哪个包含关系或逼近结论能够发挥作用。

\subsection{淡收敛}
\label{b1:subsec:120102}

\begin{definition}\label{b1:def:vague-convergence}
设 \(X\) 为局部紧 Hausdorff 空间，\(\mu_n,\mu\) 为其上的 Radon 测度。
若对每个 \(f\in C_c(X)\) 都有
\[
 \lim_{n\to\infty}\int_X f\,\mathrm{d}\mu_n=\int_X f\,\mathrm{d}\mu,
\]
则称 \(\mu_n\) \term[danshoulian]{淡收敛}[vague convergence]于 \(\mu\)，
记为 \(\mu_n\xrightarrow{v}\mu\)。
\end{definition}
\symidx{\(\mu_n\xrightarrow{v}\mu\)：由连续紧支集函数测试的淡收敛}

中文名称采用严加安的《测度论讲义》\cite[第6.5节]{Yan2004Cedulun}。
不同文献对 vague 的测试空间并不一致；例如 Fremlin 的《Measure Theory》%
\cite[437J(c)]{Fremlin2016Measure}使用 \(C_b(X)\) 定义该名称。
本书始终按上面的 \(C_c(X)\) 约定，不凭相同的外文名称移用结论。

\ProseParagraphBreak

淡极限若存在便唯一：两个候选极限在全部 \(C_c(X)\) 上有相同积分，
由正测度表示的唯一性即相等。它还自动控制每一块固定紧集上的质量。

\begin{proposition}\label{b1:prop:vague-local-bound}
若 \(\mu_n\xrightarrow{v}\mu\)，则对每个紧集 \(K\subset X\)，
有 \(\sup_n\mu_n(K)<\infty\)。
\end{proposition}
\begin{proof}
由\cref{b1:lem:lch-cutoff}取 \(h\in C_c(X)\)，使 \(0\leq h\leq1\)、
\(h|_K=1\)。于是
\[
 \mu_n(K)\leq\int_X h\,\mathrm{d}\mu_n\longrightarrow\int_X h\,\mathrm{d}\mu<\infty.
\]
收敛数列有界，结论随之成立。
\end{proof}

此处的界依赖 \(K\)，不能直接对所有紧集取上确界。例如在 \(\mathbb R\) 上令
\(\mu_n=n\delta_n\)。每个紧支集测试函数最终在点 \(n\) 处为零，
所以 \(\mu_n\xrightarrow{v}0\)，但 \(\mu_n(\mathbb R)=n\)。
甚至 \(C_0(\mathbb R)\) 中的函数 \(f(x)=(1+|x|)^{-1/2}\) 也满足
\(\int f\,\mathrm{d}\mu_n=n/\sqrt{1+n}\to\infty\)。
因此没有整体界时，不能把紧支集测试悄悄改为无穷远消失的测试。

\begin{proposition}\label{b1:prop:vague-c0-bounded}
设 \(\mu_n,\mu\) 为局部紧 Hausdorff 空间上的有限正 Radon 测度，并且
\(M=\sup_n\mu_n(X)<\infty\)。则 \(\mu_n\xrightarrow{v}\mu\) 当且仅当
它们在 \(C_0(X)^*\) 中弱*\WeakStarJoin{}收敛于 \(\mu\)。
对有限符号或复 Radon 测度，若改用 \(\sup_n|\mu_n|(X)<\infty\)，结论仍成立。
\end{proposition}
\begin{proof}
弱*\WeakStarJoin{}收敛蕴含淡收敛，因为 \(C_c(X)\subset C_0(X)\)。反过来，给定
\(f\in C_0(X)\) 和 \(\varepsilon>0\)，由\cref{b1:prop:cc-dense-c0}%
选 \(g\in C_c(X)\)，使 \(\|f-g\|_\infty<\varepsilon\)。于是
\[
 \left|\int_X f\,\mathrm{d}\mu_n-\int_X f\,\mathrm{d}\mu\right|
 \leq\left|\int_X g\,\mathrm{d}\mu_n-\int_X g\,\mathrm{d}\mu\right|
       +\varepsilon\bigl(M+\mu(X)\bigr).
\]
先令 \(n\to\infty\)，再令 \(\varepsilon\downarrow0\)，便得到所需收敛。
一般标量测度的证明相同，把两个质量替换为全变差质量即可。
\end{proof}

上式中的统一界有明确用途：同一个一致逼近误差，必须对所有 \(n\) 同时有效。
前一例正是在这一步失去控制。另一方面，若已知在 \(C_0(X)^*\) 中弱*\WeakStarJoin{}收敛，
一致有界原理又保证范数一致有界，因此对于收敛的序列，弱*\WeakStarJoin{}条件本身包含了全变差界。

\subsection{弱收敛}
\label{b1:subsec:120103}

\begin{definition}\label{b1:def:measure-weak-convergence}
设 \(X\) 为拓扑空间，\(\mu_n,\mu\) 为其上的有限正 Borel 测度。
若对每个 \(f\in C_b(X)\) 都有
\[
 \lim_{n\to\infty}\int_X f\,\mathrm{d}\mu_n=\int_X f\,\mathrm{d}\mu,
\]
则称 \(\mu_n\) \term[cedu-ruoshoulian]{弱收敛}[weak convergence of measures]于 \(\mu\)，
记为 \(\mu_n\Rightarrow\mu\)。测度空间上使所有映射
\(\nu\mapsto\int f\,\mathrm{d}\nu\)、\(f\in C_b(X)\) 连续的最弱拓扑，称为这里的弱拓扑。
\end{definition}
\symidx{\(\mu_n\Rightarrow\mu\)：由有界连续函数测试的测度弱收敛}

这个拓扑在 \(\mu\) 处的基本邻域可写成
\[
 \left\{\nu:
 \left|\int_X f_j\,\mathrm{d}\nu-\int_X f_j\,\mathrm{d}\mu\right|<\varepsilon,
 \quad 1\leq j\leq m\right\},
\]
其中只选有限多个测试函数。对于一般拓扑空间，连续函数未必足以区别测度；
后文主要使用度量空间，或局部紧空间上的 Radon 测度，这两个场合没有这样的歧义。

\begin{proposition}\label{b1:prop:weak-mass-uniqueness}
若 \(\mu_n\Rightarrow\mu\)，则 \(\mu_n(X)\to\mu(X)\)。
在度量空间上，有限正 Borel 测度由其在 \(C_b(X)\) 上的积分唯一确定；
在局部紧 Hausdorff 空间上，有限正 Radon 测度也具有此性质。
\end{proposition}
\begin{proof}
第一点取 \(f=1\)。在度量空间上，对非空闭集 \(F\) 定义
\[
 h_k(x)=\max\{1-k\cdot \mathrm{d}(x,F),0\}.
\]
这些函数有界连续，且 \(h_k\downarrow\mathbf1_F\)。若两个有限测度的连续测试积分相同，
由支配收敛定理，它们在 \(F\) 上的值也相同。空集的值当然相同。
所有闭集构成包含 \(X\) 的 \(\pi\)-系并生成 Borel \(\sigma\)-代数，
故有限测度的唯一性定理给出两测度相同。局部紧情形则直接使用
\(C_c(X)\subset C_b(X)\) 及前一章的 Radon 表示唯一性。
\end{proof}

这里的测度弱收敛，不是把测度 Banach 空间 \(M(X)=C_0(X)^*\) 的整个对偶再拿来测试。
后者对应 \(\sigma\bigl(M(X),M(X)^*\bigr)\)，是\chapref{b1:ch:09}意义下 Banach 空间的弱拓扑。
本章固定 \(C_b(X)\) 测试，不能仅因名称中都有“弱”字便将它们等同。
即使 \(X\) 紧，测度弱拓扑对应的也是 \(C(X)^*\) 的弱*\WeakStarJoin{}拓扑，而非一般的弱拓扑。

\subsection{概率测度情形}
\label{b1:subsec:120104}

记 \(\mathcal P(X)\) 为 \(X\) 上所有 Borel 概率测度构成的集合，赋予刚定义的弱拓扑。
\symidx{\(\mathcal P(X)\)：Borel 概率测度空间}
若有限正测度列的每一项都有质量一，那么任何弱极限仍有质量一。
但淡极限不一定如此；缺少的质量可能已经离开每一个固定的紧区域。
\secref{b1:sec:1203}将用紧族把这种现象精确表述出来。

\ProseParagraphBreak

先看与质量逃逸不同的集中现象。在 \(\mathbb R\) 上令
\(\mathrm{d}\mu_n(x)=n\cdot\mathbf1_{[0,1/n]}(x)\,\mathrm{d}x\)。对任意 \(f\in C_b(\mathbb R)\)，有
\[
 \left|\int_{\mathbb R}f\,\mathrm{d}\mu_n-f(0)\right|
 \leq\sup_{0\leq x\leq1/n}|f(x)-f(0)|\longrightarrow0.
\]
所以 \(\mu_n\Rightarrow\delta_0\)。每个 \(\mu_n\) 都由密度给出，极限却是点质量。
此外，\(\mu_n(\{0\})=0\) 而 \(\delta_0(\{0\})=1\)，说明弱收敛不要求逐个
Borel 集的测度值收敛。下一节会指出哪些集合仍能保留这一结论。

\begin{proposition}[连续映射原理]\label{b1:prop:weak-continuous-mapping}
设 \(T:X\to Y\) 连续，\(\mu_n,\mu\) 为 \(X\) 上的有限正 Borel 测度。
若 \(\mu_n\Rightarrow\mu\)，则推前测度满足 \(T_\#\mu_n\Rightarrow T_\#\mu\)。
事实上，\(\mu\mapsto T_\#\mu\) 关于两边的弱拓扑连续。
\end{proposition}
\begin{proof}
对 \(g\in C_b(Y)\)，复合函数 \(g\circ T\) 属于 \(C_b(X)\)。由推前积分公式，
\[
 \int_Y g\,\mathrm{d}(T_\#\mu_n)=\int_X g\circ T\,\mathrm{d}\mu_n
 \longrightarrow\int_X g\circ T\,\mathrm{d}\mu=\int_Y g\,\mathrm{d}(T_\#\mu).
\]
拓扑连续性也由同一等式给出：目标空间中每个测试积分与推前映射的复合，
都是源空间中的连续测试积分。
\end{proof}

称拓扑空间之间的连续映射 \(T:X\to Y\) 为\term[zhen-yingshe]{真映射}[proper map]，
若 \(Y\) 中每个紧集的原像都是 \(X\) 中的紧集。上面的连续映射原理不要求
\(T\) 具有这一性质，因为有界连续函数的复合仍有界连续；若改用 \(C_c(Y)\) 中的
紧支集测试，则 \(g\circ T\) 的支集一般未必紧。因而用淡收敛讨论推前测度时，
通常需要真映射或其他能够控制复合函数支集的条件；\chapref{b1:ch:11}的相应习题采用的正是这一条件。
在概率论中，若随机变量 \(Z_n\) 的分布为 \(\mu_n\)，它们分布的弱收敛也称为%
\term[yifenbu-shoulian]{依分布收敛}[convergence in distribution]，通常记作
\(Z_n\xrightarrow{d}Z\)，其中 \(Z\) 的分布为 \(\mu\)。
这一说法只比较分布，不要求随机变量处在同一个概率空间上。

\ProseParagraphBreak

后面的概率论版本允许 \(X\) 为无限维度量空间。在无限维 Hilbert 空间中，
连续紧支集函数可能只有零函数，\chapref{b1:ch:11}已说明其原因。
因此不能先用 \(C_c(X)\) 定义所有概率测度的收敛，再期待局部紧情形的论证自动推广。
而 \(C_b(X)\) 测试则在这一情形仍有充分的区分能力。

\ProseParagraphBreak

\cref{b1:tab:vague-weak-measures}概括两种收敛最直接的分界。淡收敛只测试局部可见的质量；弱收敛还用常函数 \(1\) 测试总质量，因此概率测度的弱极限不会丢失质量。
\begin{table}[htbp]
  \centering
  \caption[淡收敛与弱收敛]{淡收敛与弱收敛。表中淡收敛一行以局部紧 Hausdorff 空间上的 Radon 测度为背景。}
  \label{b1:tab:vague-weak-measures}
  \begin{tabularx}{\textwidth}{@{}p{2.4cm}p{4.0cm}X@{}}
    \toprule
    收敛方式 & 测试函数 & 对逃向无穷远质量的反应 \\
    \midrule
    淡收敛 & \(C_c(X)\) & 只观察固定紧区域，因而可能看不见逃逸质量 \\
    弱收敛 & \(C_b(X)\) & 常函数 \(1\) 同时控制总质量；概率极限仍是概率测度 \\
    \bottomrule
  \end{tabularx}
\end{table}

\cref{b1:fig:visual-f07}以原子位置说明测试类与紧族条件的作用。
% F07：本书原创矢量示意；依据所在节的定义、证明或明确模型。
\begin{figure}[htbp]
\centering
\begin{tikzpicture}[x=1cm,y=1cm,>=stealth,every node/.style={font=\small},line width=.65pt]

\draw[->,gray](0,0)--(5.7,0)node[right]{\(x\)};
\foreach \x/\n in {3/1,1.5/2,.75/4}{\draw[->,very thick,AnalysisBlue](\x,0)--(\x,1.8);\node[above]at(\x,1.8){\(1/\n\)};}
\draw[->,orange!80!black](2.8,2.6)--(.3,2.6);
\node at(2.5,3.25){\(\delta_{1/n}\Rightarrow\delta_0\)};
\draw[->,gray](7,0)--(12.7,0)node[right]{\(x\)};
\fill[AnalysisBlue!10](7,0)rectangle(9.3,1.4);
\node at(8.15,.65){固定紧区间};
\foreach \x in {8,10,12}{\draw[->,very thick,AnalysisBlue](\x,0)--(\x,1.8);}
\draw[->,orange!80!black](9,2.6)--(12.5,2.6);
\node at(9.8,3.25){\(\delta_n\) 淡收敛到零};
\node at(9.8,-.65){总质量始终为 \(1\)};

\end{tikzpicture}
\caption[原子测度的汇聚与质量逃逸]{原子测度的汇聚与质量逃逸。每根箭头都代表单位质量。逃逸族的紧支集测试趋零，但常函数一的积分不变，因此它不弱收敛到零测度。}
\label{b1:fig:visual-f07}
\end{figure}

